% gallery-deck.tex — BayesStack Academic Editorial System (Gallery 3.0)
%
% PURPOSE: Demonstrates 30 Essential Visual Teaching Compositions for world-class
%   computational, mathematical, and theoretical university-grade slide decks.
\documentclass[aspectratio=169,10pt]{beamer}
\usepackage{import}
\usepackage[science]{bayesstackslides}

\CourseMetadata{GALLERY 3.0}{Bayes Academic Editorial System}{30 Visual Teaching Compositions}
\author{BayesStack Academic Group}
\date{}

\begin{document}

% =============================================================================
% 00. TITLE SLIDE (Academic Gravitas)
% =============================================================================
\BayesCourseTitleSlide{Bayes Academic Editorial System}{30 Essential Compositions for World-Class Computational \& Theoretical Instruction}

% =============================================================================
% 01. SINGLE BIG IDEA (Monumental Scale Contrast)
% =============================================================================
\begin{frame}{01 — Single Big Idea}
  \BigIdeaHero{%
    Probability is the logic of science under uncertainty.%
  }{%
    Parameters are random variables quantifying states of knowledge, not asymptotic physical frequencies.%
  }{%
    Evidence reweights a prior belief to form a posterior through Bayes' Theorem.%
  }
\end{frame}

% =============================================================================
% 02. BAYES' THEOREM GEOMETRIC PROBABILITY REWEIGHTING
% =============================================================================
\begin{frame}{02 — Probability Mass Geometric Reweighting}
  \FigureDominant[.62]{%
    \begin{tikzpicture}
      \begin{axis}[
        bayes/science plot,
        xlabel={\small Parameter Hypothesis $\theta$},
        ylabel={\small Density Mass},
        domain=0:1,
        samples=100,
        xmin=0, xmax=1, ymin=0, ymax=4.6,
        legend pos=north west,
        width=.98\linewidth, height=.54\textheight
      ]
        \addplot[thick, color=BayesSecondary, dashed] {1};
        \addlegendentry{\scriptsize Prior $p(\theta) = \text{Beta}(1,1)$}
        
        \addplot[thick, color=BayesAmber] {3.2 * exp(-0.5*((x-0.78)/0.18)^2)};
        \addlegendentry{\scriptsize Likelihood $p(\mathcal{D} \mid \theta)$}
        
        \addplot[very thick, color=BayesTeal, fill=BayesTeal!12!white] {4.2 * (x^7) * ((1-x)^2) * 36 / 3.36} \closedcycle;
        \addlegendentry{\scriptsize Posterior $p(\theta \mid \mathcal{D}) \propto p(\mathcal{D}\mid\theta)p(\theta)$}
      \end{axis}
    \end{tikzpicture}%
  }{The Reweighting Principle}{%
    \begin{itemize}\setlength{\itemsep}{.25em}
      \item \textbf{Prior $p(\theta)$}: Baseline hypothesis density before evidence.
      \item \textbf{Likelihood $p(\mathcal{D}\mid\theta)$}: Multiplicative filter derived from empirical observations.
      \item \textbf{Posterior $p(\theta\mid\mathcal{D})$}: Renormalized density surviving the likelihood filter.
    \end{itemize}%
  }
  \vfill
  \BayesCaption{Geometric Intuition: Observation acts as an analytical filter contracting probability mass around empirical truth.}
\end{frame}

% =============================================================================
% 03. FIGURE-DOMINANT CANVAS: CONTOUR LOSS LANDSCAPE
% =============================================================================
\begin{frame}{03 — Ill-Conditioned Convex Optimization}
  \FigureDominant[.64]{%
    \begin{tikzpicture}[scale=0.92]
      % Contour ellipses
      \draw[BayesBorder!80!BayesDark,thick] (0,0) ellipse (3.6cm and 2.1cm);
      \draw[BayesBorder!90!BayesDark,thick] (0,0) ellipse (2.7cm and 1.5cm);
      \draw[BayesBorder!95!BayesDark,thick] (0,0) ellipse (1.8cm and 0.95cm);
      \draw[BayesTeal!60!white,thick] (0,0) ellipse (0.9cm and 0.45cm);
      \fill[BayesTeal] (0,0) circle (2.2pt) node[below=3pt,font=\scriptsize\bfseries\color{BayesTeal}] {Global Minimum $\theta^*$};
      
      % High-curvature direction annotation
      \fill[BayesAmber] (-2.2,0.8) circle (2.0pt) node[above=2pt,font=\tiny\bfseries\color{BayesAmber}] {Steep Curvature};
      
      % SGD Path (Noisy zig-zag)
      \draw[thick,color=BayesCrimson,-{Latex[length=1.6mm]}] (-3.2,1.8) -- (-2.6,1.1) -- (-2.4,1.4) -- (-1.7,0.7) -- (-1.5,1.0) -- (-0.8,0.4) -- (-0.3,0.5) -- (-0.05,0.08);
      \node[font=\tiny\bfseries\color{BayesCrimson},above left] at (-3.2,1.8) {SGD Trajectory};
      
      % Momentum Path (Smooth accelerated curve)
      \draw[very thick,color=BayesBlue,-{Latex[length=2.0mm]}] (-3.2,1.5) to[out=-20,in=160] (-1.2,0.25) to[out=-20,in=170] (-0.02,0.02);
      \node[font=\tiny\bfseries\color{BayesBlue},below left] at (-3.2,1.5) {Momentum SGD};
      
      % Gradient Vector
      \draw[thick,color=BayesDark,-{Latex[length=2.0mm]}] (1.8,0.95) -- (1.1,0.58) node[midway,above right,font=\tiny\color{BayesDark}] {$-\nabla \mathcal{L}$};
    \end{tikzpicture}%
  }{Ill-Conditioned Curvature}{%
    \begin{itemize}\setlength{\itemsep}{.25em}
      \item \textbf{Ravine Anisotropy}: High condition number $\kappa = \lambda_{\max}/\lambda_{\min} \gg 1$ causes standard SGD to oscillate orthogonally.
      \item \textbf{Momentum Accumulation}: Velocity $v_t = \beta v_{t-1} + \eta \nabla \mathcal{L}$ cancels transverse oscillations while accelerating down the valley.
    \end{itemize}%
  }
  \vfill
  \BayesCaption{Canonical convex geometry: path comparison on an anisotropic quadratic loss surface.}
\end{frame}

% =============================================================================
% 04. CHALKBOARD HERO: GAUSSIAN LOG-LIKELIHOOD
% =============================================================================
\begin{frame}{04 — Monumental Chalkboard Equation}
  \vfill
  \ChalkboardHero{%
    \ln p(\mathbf{y} \mid \mathbf{X}, \mathbf{w}, \sigma^2) = \underbrace{-\frac{N}{2}\ln(2\pi\sigma^2)}_{\text{Normalization Constant}} - \underbrace{\frac{1}{2\sigma^2} \|\mathbf{y} - \mathbf{X}\mathbf{w}\|_2^2}_{\text{Residual Sum of Squares}}%
  }{%
    Log-Likelihood Decomposition: Maximizing Gaussian log-likelihood is mathematically identical to minimizing quadratic Euclidean error.%
  }
  \vfill
  \TakeawayBanner{Least-squares regression is the exact Maximum Likelihood Estimator under homoscedastic Gaussian noise.}
\end{frame}

% =============================================================================
% 05. MULTI-STEP MATHEMATICAL DERIVATION (Blackboard Flow)
% =============================================================================
\begin{frame}{05 — Derivation Sequence (Blackboard Clarity)}
  \AnnotatedStep%
    {1}%
    {Likelihood $\times$ Prior Kernel}%
    {p(\theta \mid y) \propto \left[\theta^y (1-\theta)^{n-y}\right] \cdot \left[\theta^{\alpha-1}(1-\theta)^{\beta-1}\right]}%
    {Direct product of the Binomial observation density and the conjugate Beta prior kernel, dropping normalising constants independent of $\theta$.}%
  
  \vspace{.35cm}
  \AnnotatedStep%
    {2}%
    {Hyperparameter Addition}%
    {p(\theta \mid y) \propto \theta^{(\alpha + y) - 1}\,(1-\theta)^{(\beta + n - y) - 1} \implies \mathrm{Beta}(\alpha + y, \, \beta + n - y)}%
    {Recognise the updated kernel directly: Prior pseudo-counts $\alpha, \beta$ update additively with zero numerical integration error.}
  
  \vfill
  \BayesCaption{Classroom Standard: Mathematical expressions take 60\% of canvas; annotations remain clearly subordinate.}
\end{frame}

% =============================================================================
% 06. THEOREM & PROOF SKETCH
% =============================================================================
\begin{frame}{06 — Formal Theorem \& Proof Sketch}
  \begin{TheoremBox}{Bernstein--von Mises Theorem (Asymptotic Posterior Normality)}
    Let $X_1, \dots, X_n \overset{\text{i.i.d.}}{\sim} p(x \mid \theta_0)$. Under smooth identifiability conditions, as sample size $n \to \infty$:
    \[
      \left\| p(\theta \mid X_{1:n}) - \mathcal{N}\left(\hat{\theta}_{\mathrm{MLE}}, \, [n\mathcal{I}(\theta_0)]^{-1}\right) \right\|_{\mathrm{TV}} \xrightarrow{\mathrm{a.s.}} 0
    \]
  \end{TheoremBox}
  \vspace{.10cm}
  \begin{ProofSketch}{Taylor Expansion of Log-Posterior}
    Expand $\ln p(\theta \mid \mathcal{D}) = \ln p(\hat{\theta}) + \frac{1}{2}(\theta - \hat{\theta})^T \nabla^2 \ln p(\hat{\theta}) (\theta - \hat{\theta}) + \mathcal{O}(\|\theta - \hat{\theta}\|^3)$. By the Law of Large Numbers, $-\frac{1}{n}\nabla^2 \ln p(\hat{\theta}) \to \mathcal{I}(\theta_0)$, while the prior $p(\theta)$ contributes an asymptotically vanishing $\mathcal{O}(1/n)$ term. \hfill$\blacksquare$
  \end{ProofSketch}
\end{frame}

% =============================================================================
% 07. EDITORIAL SPLIT NARRATIVE (Hairline Separation)
% =============================================================================
\begin{frame}{07 — Editorial Split Narrative}
  \SplitEditorial%
    {Frequentist Paradigm}{%
      \begin{itemize}\setlength{\itemsep}{.25em}
        \item Parameters $\theta$ are fixed, immutable constants in nature.
        \item Probabilities denote long-run limiting frequencies under infinite hypothetical repetitions.
        \item Performance is measured across the data sampling distribution (e.g. unbiasedness, coverage).
      \end{itemize}
    }%
    {Bayesian Paradigm}{%
      \begin{itemize}\setlength{\itemsep}{.25em}
        \item Parameters $\theta$ are random variables encoding epistemic uncertainty.
        \item Probabilities denote subjective degrees of rational belief conditioned on observed evidence.
        \item All conclusions derive from the single realized posterior distribution $p(\theta \mid \mathcal{D})$.
      \end{itemize}
    }
  \vfill
  \BayesCaption{Pedagogical Intent: Direct conceptual contrast without unnecessary card containers.}
\end{frame}

% =============================================================================
% 08. CODE WALKTHROUGH & LINE HIGHLIGHTS
% =============================================================================
\begin{frame}[fragile]{08 — Production Probabilistic Kernel}
  \begin{columns}[T,onlytextwidth]
    \begin{column}{.60\textwidth}
\begin{HighlightedCodeBlock}[numbers=left]{Python}{7,8,9,10}
import numpy as np

def metropolis_hastings(target_log_prob, init,
                        n_samples=2000):
    samples = [init]
    for _ in range(n_samples):
        current = samples[-1]
        proposal = current + np.random.normal(0, 0.45)
        log_alpha = target_log_prob(proposal) - target_log_prob(current)
        if (np.log(np.random.uniform())
                < log_alpha):
            samples.append(proposal)
        else:
            samples.append(current)
    return np.array(samples)
\end{HighlightedCodeBlock}
    \end{column}
    \hfill
    \begin{column}{.37\textwidth}
      \vspace{.10cm}
      {\small\bfseries\color{BayesTeal}Metropolis Acceptance Ratio}\par\vspace{.08cm}
      {\small Log-space evaluation prevents floating point underflow:}\par\vspace{.08cm}
      \[
        \alpha = \min\left(1, \, \frac{p(\theta^*)}{p(\theta^{(t)})}\right)
      \]
      \vspace{.08cm}
      \BayesQuietNote{\textbf{Random-walk scaling heuristic:} Under the usual asymptotic assumptions, targets near $44\%$ in 1D and $23.4\%$ in high dimension are often efficient---they are not universal rates.}
    \end{column}
  \end{columns}
\end{frame}

% =============================================================================
% 09. LIVE RUNTIME EXECUTION & TELEMETRY
% =============================================================================
\begin{frame}[fragile]{09 — Code + Interactive Execution Telemetry}
  \begin{columns}[T,onlytextwidth]
    \begin{column}{.52\textwidth}
\begin{CodeBlock}{Python}
import jax.numpy as jnp
from jax import grad, jit, vmap

@jit
def loss_fn(w, x, y):
    preds = jnp.dot(x, w)
    return jnp.mean((preds - y)**2)

grad_fn = jit(grad(loss_fn))
\end{CodeBlock}
      \vspace{.08cm}
      \BayesQuietNote{JAX automatically fuses computation kernels and generates exact analytic gradients via reverse-mode AD.}
    \end{column}
    \hfill
    \begin{column}{.45\textwidth}
      \TerminalWindow[GPU Telemetry · NVIDIA A100]{%
        >>> compiled\_fn = grad\_fn.lower(w, x, y).compile() \\
        >>> timeit(compiled\_fn) \\
        3.12 $\mu$s $\pm$ 41 ns per loop (XLA fused) \\
        >>> max\_memory\_allocated() \\
        12.4 MB (Zero host-device copy overhead)
      }
      \vspace{.08cm}
      \DataProvenance[COMPUTE BENCHMARK]{JAX 0.4.28 · CUDA 12.4}
    \end{column}
  \end{columns}
\end{frame}

% =============================================================================
% 10. ALGORITHMIC SPECIFICATION (Cognitive State Trace)
% =============================================================================
\begin{frame}{10 — Algorithmic Specification: SGHMC}
  \begin{tcolorbox}[enhanced,arc=1.2mm,boxrule=.4pt,colback=white,colframe=BayesBorder,borderline west={2.2pt}{0pt}{BayesTeal},title={Algorithm: Stochastic Gradient Hamiltonian Monte Carlo},coltitle=BayesTeal,fonttitle=\bfseries]
    \begin{enumerate}\setlength{\itemsep}{.08cm}\small
      \item \textbf{Initialize} position $\theta_0 \sim p_0$, momentum $v_0 \sim \mathcal{N}(0, M)$, friction parameter $\gamma$.
      \item \textbf{For} iteration $t = 1, \dots, T$:
        \begin{itemize}\footnotesize
          \item Sample stochastic mini-batch gradient $\tilde{g}_t(\theta_t) \approx -\nabla_\theta \log p(\theta_t \mid \mathcal{D})$.
          \item Update momentum with noise injection: $v_{t+1} \leftarrow v_t - \epsilon \tilde{g}_t(\theta_t) - \gamma v_t + \mathcal{N}(0, 2\gamma \epsilon M)$.
          \item Update parameter coordinate: $\theta_{t+1} \leftarrow \theta_t + \epsilon M^{-1} v_{t+1}$.
        \end{itemize}
      \item \textbf{Output} stationary ergodic Markov chain sample collection $\{\theta_1, \dots, \theta_T\}$.
    \end{enumerate}
  \end{tcolorbox}
  \vfill
  \BayesCaption{Computational Complexity: $\mathcal{O}(B)$ per leapfrog step, where $B \ll N$ is the mini-batch dimension.}
\end{frame}

% =============================================================================
% 11. DISTRIBUTED ARCHITECTURE (Uncluttered Topology)
% =============================================================================
\begin{frame}{11 — Distributed Parameter Server Topology}
  \centering
  \vspace{.15cm}
  \begin{tikzpicture}[
      box/.style={bayes/node,minimum width=2.4cm,minimum height=.74cm},
      lbl/.style={font=\tiny\bfseries,inner sep=2pt}
    ]
    % Parameter Server (Top Center)
    \node[bayes/accent,fill=BayesDark,minimum width=3.4cm,minimum height=.8cm] (ps) at (0, 2.2) {\small\color{white} Parameter Server};
    
    % Workers (Bottom Row)
    \node[bayes/model,box] (w1) at (-3.8, 0) {\small Worker 01};
    \node[bayes/model,box] (w2) at (0, 0) {\small Worker 02};
    \node[bayes/model,box] (w3) at (3.8, 0) {\small Worker 03};
    
    % Upward Flow: Gradients (Push)
    \draw[bayes/flow,color=BayesCrimson] (w1.north) to[bend left=12] node[lbl,color=BayesCrimson,above left] {Push $\nabla \mathcal{L}_1$} (ps.west);
    \draw[bayes/flow,color=BayesCrimson] (w2.north) -- node[lbl,color=BayesCrimson,left=2pt] {Push $\nabla \mathcal{L}_2$} (ps.south);
    \draw[bayes/flow,color=BayesCrimson] (w3.north) to[bend right=12] node[lbl,color=BayesCrimson,above right] {Push $\nabla \mathcal{L}_3$} (ps.east);
    
    % Downward Flow: Parameters (Pull)
    \draw[bayes/flow,dashed,color=BayesTeal] (ps.south west) to[bend right=16] node[lbl,color=BayesTeal,below left] {Pull $\mathbf{w}$} (w1.north east);
    \draw[bayes/flow,dashed,color=BayesTeal] (ps.south east) to[bend left=16] node[lbl,color=BayesTeal,below right] {Pull $\mathbf{w}$} (w3.north west);
  \end{tikzpicture}
  \vfill
  \BayesCaption{Architecture Invariant: Workers evaluate local mini-batches in parallel; asynchronous updates eliminate barrier synchronization.}
\end{frame}

% =============================================================================
% 12. EMPIRICAL BENCHMARK TABLE & PROVENANCE
% =============================================================================
\begin{frame}{12 — Empirical Benchmark Evaluation}
  \begin{BayesTable}{lcccc}
    \BayesTableHeader
    \BayesTableHead{Model} & \BayesTableHead{Params} & \BayesTableHead{Latency} & \BayesTableHead{Top-1} & \BayesTableHead{Memory} \\
    ResNet-50 & 25.6M & 4.2\,ms & 76.1\% & 1.2\,GB \\
    Vision Transformer (ViT-B/16) & 86.0M & 9.8\,ms & 81.8\% & 3.4\,GB \\
    Swin Transformer (Swin-L) & 197.0M & 21.4\,ms & 86.3\% & 7.8\,GB \\
    BayesStack Diffusion Backbone & 310.0M & 34.2\,ms & \textbf{89.4\%} & 11.2\,GB \\
  \end{BayesTable}
  \vspace{.20cm}
  \begin{columns}[c,onlytextwidth]
    \begin{column}{.54\textwidth}
      {\small\bfseries\color{BayesInk}Inference Scaling Performance}\par\vspace{.04cm}
      \BayesQuietNote{Evaluated with TensorRT fp16 compilation across 50,000 ImageNet validation images.}
    \end{column}
    \hfill
    \begin{column}{.44\textwidth}
      \raggedleft
      \DataProvenance[SYNTHETIC BENCHMARK]{B32 · TRT 10.0}
    \end{column}
  \end{columns}
\end{frame}

% =============================================================================
% 13. MISCONCEPTION CORRECTION & RESOLUTION
% =============================================================================
\begin{frame}{13 — Misconception Correction}
  \begin{columns}[T,onlytextwidth]
    \begin{column}{.48\textwidth}
      {\small\bfseries\color{BayesCrimson}Fallacy: Flat Priors Are Always Uninformative}\par\BayesGap
      {\small A uniform prior $p(\theta) \propto 1$ is \textbf{not reparameterization-invariant}. A flat prior on variance $\sigma^2$ implies an aggressively informative prior on standard deviation $\sigma$ or precision $\tau = 1/\sigma^2$.}\par\BayesGap
      \BayesQuietNote{Transformations distort Euclidean density volume elements by the Jacobian determinant $|\frac{d\theta}{d\phi}|$.}
    \end{column}
    \hfill{\color{BayesBorder}\vrule width .4pt}\hfill
    \begin{column}{.48\textwidth}
      {\small\bfseries\color{BayesTeal}Resolution: Jeffreys' Invariant Prior}\par\BayesGap
      \[
        p_{\text{Jeffreys}}(\theta) \propto \sqrt{\det \mathcal{I}(\theta)}
      \]
      \BayesGap
      {\small Under any smooth invertible mapping $\phi = g(\theta)$, the transformed prior satisfies $p(\phi) = p(\theta) |\frac{d\theta}{d\phi}|$.}\par\BayesGap
      \BayesQuietNote{Fisher information transforms as a 2-tensor on a Riemannian statistical manifold.}
    \end{column}
  \end{columns}
  \vfill
  \BayesCaption{Pedagogical Lesson: Objectivity in prior specification requires geometric invariance, not naive flatness.}
\end{frame}

% =============================================================================
% 14. WORKED QUANTITATIVE EXAMPLE (Black-Scholes Call)
% =============================================================================
\begin{frame}{14 — Worked Quantitative Derivation}
  \begin{columns}[T,onlytextwidth]
    \begin{column}{.48\textwidth}
      {\small\bfseries\color{BayesInk}Problem Setup}\par\BayesSmallGap
      {\small Price European call with spot $S_0 = 100$, strike $K = 100$, rate $r = 5\%$, vol $\sigma = 20\%$, maturity $T = 1\text{ yr}$.}\par\BayesGap
      {\small\bfseries\color{BayesTeal}Step 1: Compute $d_1$ and $d_2$}\par\BayesGap
      \[
        d_1 = \frac{\ln(100/100) + (0.05 + \frac{0.04}{2})\cdot 1}{0.20 \cdot 1} = 0.35
      \]
      \BayesSmallGap
      \[
        d_2 = d_1 - \sigma\sqrt{T} = 0.35 - 0.20 = 0.15
      \]
    \end{column}
    \hfill
    \begin{column}{.48\textwidth}
      {\small\bfseries\color{BayesTeal}Step 2: Normal CDF Evaluation}\par\BayesGap
      \[
        N(d_1) = 0.6368, \quad N(d_2) = 0.5596
      \]
      \BayesGap
      {\small\bfseries\color{BayesDark}Step 3: Analytical Call Option Value}\par\BayesGap
      \[
        C = 100(0.6368) - 100e^{-0.05}(0.5596) = \mathbf{\$10.45}
      \]
      \BayesGap
      \BayesQuietNote{\textbf{Interpretation:} Delta hedge ratio $\Delta = N(d_1) = 0.6368$ shares per contract.}
    \end{column}
  \end{columns}
  \vfill
  \BayesCaption{Classroom Standard: Full arithmetic clarity with intermediate values displayed.}
\end{frame}

% =============================================================================
% 15. SOCRATIC QUESTION & INTUITION PROBE
% =============================================================================
\begin{frame}{15 — Socratic Inquiry: Regularization Priors}
  \vfill
  \begin{center}
    {\fontsize{14}{18}\selectfont\bfseries\color{BayesInk}Why does L2 weight decay produce dense shrinkage, while L1 regularization produces exact parameter sparsity?}
  \end{center}
  \vspace{.25cm}
  \begin{columns}[T,onlytextwidth]
    \begin{column}{.48\textwidth}
      {\small\bfseries\color{BayesBlue}L2 Decay $\Longleftrightarrow$ Gaussian Prior $\mathcal{N}(0, \sigma^2)$}\par\vspace{.06cm}
      {\small Quadratic log-penalty $-\frac{w^2}{2\sigma^2}$ has derivative $\nabla \propto w$, smoothly vanishing as $w \to 0$ without forcing weights to zero.}
    \end{column}
    \hfill
    \begin{column}{.48\textwidth}
      {\small\bfseries\color{BayesCrimson}L1 Penalty $\Longleftrightarrow$ Laplace Prior $\text{Laplace}(0, b)$}\par\vspace{.06cm}
      {\small Linear log-penalty $-\frac{|w|}{b}$ maintains constant non-zero subgradient $\pm 1/b$ at the origin, driving non-informative weights to exact zero.}
    \end{column}
  \end{columns}
  \vfill
\end{frame}

% =============================================================================
% 16. ACTIVE LAB CHALLENGE & CODE TARGET
% =============================================================================
\begin{frame}{16 — Active Hands-on Lab Challenge}
  \begin{tcolorbox}[enhanced,arc=1.2mm,boxrule=.4pt,colback=BayesPale!50!white,colframe=BayesBorder,borderline west={2.2pt}{0pt}{BayesTeal},title={Lab 04: Vectorized NUTS Sampler in JAX},coltitle=BayesDeepTeal,fonttitle=\bfseries]
    \begin{enumerate}\setlength{\itemsep}{.08cm}\small
      \item Implement the leapfrog symplectic integrator using \InlineCode{jax.vmap} across 128 parallel Markov chains.
      \item Track build-tree termination via the No-U-Turn criterion to prevent trajectory self-intersection.
      \item Compute the split-$\hat{R}$ potential scale reduction factor across warmup chains.
    \end{enumerate}
  \end{tcolorbox}
  \vfill
  \begin{columns}[c,onlytextwidth]
    \begin{column}{.53\textwidth}
      \TakeawayBanner{Target: Sample a 500-dimensional posterior in $< 800\text{ ms}$ on an A100 GPU.}
    \end{column}
    \hfill
    \begin{column}{.45\textwidth}
      \raggedleft\DataProvenance[LAB SPECIFICATION]{Autograd Evaluation Suite}
    \end{column}
  \end{columns}
\end{frame}

% =============================================================================
% 17. MULTI-STAGE PROCESS LIFECYCLE (Zero Overlap QA)
% =============================================================================
\begin{frame}{17 — MLOps Pipeline Lifecycle}
  \centering
  \vspace{.10cm}
  \begin{tikzpicture}[
      box/.style={bayes/node,minimum width=2.6cm,minimum height=.74cm},
      lbl/.style={font=\tiny,inner sep=2pt,color=BayesSecondary}
    ]
    % Row 1
    \node[box,bayes/data] (d) at (0, 1.8) {\small 1. Feature Store};
    \node[box,bayes/model] (m) at (4.6, 1.8) {\small 2. Model Training};
    \node[box,bayes/accent,fill=BayesTeal] (v) at (9.2, 1.8) {\small\bfseries\color{white} 3. Evaluation Gate};
    
    % Row 2
    \node[box,bayes/accent,fill=BayesDark] (dep) at (9.2, -0.4) {\small\bfseries\color{white} 4. Canary Deploy};
    \node[box,bayes/data] (mon) at (4.6, -0.4) {\small 5. Drift Monitor};
    \node[box,bayes/model] (ret) at (0, -0.4) {\small 6. Retrain Trigger};
    
    % Connections
    \draw[bayes/flow] (d) -- (m);
    \draw[bayes/flow] (m) -- (v);
    \draw[bayes/flow] (v) -- node[lbl,right=3pt] {Pass} (dep);
    \draw[bayes/flow] (dep) -- (mon);
    \draw[bayes/flow] (mon) -- node[lbl,above=2pt] {$\text{PSI} > 0.2$} (ret);
    \draw[bayes/flow,dashed,color=BayesTeal] (ret) -- (d);
  \end{tikzpicture}
  \vfill
  \BayesCaption{Closed-Loop Pipeline: Continuous monitoring feeds drift signals back into automated ingestion.}
\end{frame}

% =============================================================================
% 18. VISUAL COGNITIVE SEQUENCE: TRANSFORMER ATTENTION
% =============================================================================
\begin{frame}{18 — Transformer Attention: Cognitive Sequence}
  \begin{columns}[c,onlytextwidth]
    \begin{column}{.48\textwidth}
      \centering
      \begin{tikzpicture}[
          box/.style={bayes/node,minimum width=1.45cm,minimum height=.56cm,font=\scriptsize\bfseries}
        ]
        % Deliberately generous horizontal and vertical lanes: no nodes or
        % connectors share a bounding box, even after a font substitution.
        \node[box,bayes/data] (q) at (-2.0, 0) {Query $Q$};
        \node[box,bayes/data] (k) at (0, 0) {Key $K$};
        \node[box,bayes/data] (v) at (2.0, 0) {Value $V$};
        
        \node[box,bayes/model,minimum width=2.45cm] (qk) at (-1.0, -1.05) {$\text{MatMul } Q K^T$};
        \node[box,bayes/accent,fill=BayesTeal,minimum width=2.45cm] (sc) at (-1.0, -2.00) {\color{white} Scale $1/\sqrt{d_k}$};
        \node[box,bayes/accent,fill=BayesBlue,minimum width=2.45cm] (sm) at (-1.0, -2.95) {\color{white} Softmax $(\cdot)$};
        \node[box,bayes/accent,fill=BayesDark,minimum width=3.20cm] (out) at (0, -4.00) {\color{white} Value Weighted Sum};
        
        \draw[bayes/flow] (q.south) |- (qk.north);
        \draw[bayes/flow] (k.south) |- (qk.north);
        \draw[bayes/flow] (qk) -- (sc);
        \draw[bayes/flow] (sc) -- (sm);
        \draw[bayes/flow] (sm) -- (out);
        \draw[bayes/flow] (v.south) -- ++(0,-4.0) -| (out.east);
      \end{tikzpicture}
    \end{column}
    \hfill
    \begin{column}{.48\textwidth}
      {\small\bfseries\color{BayesTeal}The Cognitive Sequence}\par\BayesGap
      \begin{enumerate}\setlength{\itemsep}{.20em}\small
        \item \textbf{Match}: Compute raw token compatibility $S = QK^T$.
        \item \textbf{Scale}: Divide by $\sqrt{d_k}$ to prevent variance explosion into vanishing gradient zones.
        \item \textbf{Normalize}: Softmax maps scores to a convex probability simplex.
        \item \textbf{Aggregate}: Retrieve weighted linear combination of value vectors $V$.
      \end{enumerate}
    \end{column}
  \end{columns}
  \vfill
  \BayesCaption{Cognitive Flow: The attention formula emerges as the natural compression of a 4-step retrieval process.}
\end{frame}

% =============================================================================
% 19. QUANTITATIVE FINANCE: ITÔ'S LEMMA CHALKBOARD STEP
% =============================================================================
\begin{frame}{19 — Stochastic Calculus: Itô's Lemma}
  \AnnotatedStep%
    {1}%
    {Second-Order Taylor Expansion}%
    {df(t, X_t) = \frac{\partial f}{\partial t}dt + \frac{\partial f}{\partial x}dX_t + \frac{1}{2}\frac{\partial^2 f}{\partial x^2}(dX_t)^2}%
    {Unlike standard deterministic calculus, the quadratic variation of Brownian motion does not vanish: $(dX_t)^2 \neq 0$.}%
  
  \vspace{.35cm}
  \AnnotatedStep%
    {2}%
    {Brownian Differential Product}%
    {(dX_t)^2 = (\mu dt + \sigma dW_t)^2 = \sigma^2 dt \implies df = \left(\frac{\partial f}{\partial t} + \mu \frac{\partial f}{\partial x} + \frac{1}{2}\sigma^2 \frac{\partial^2 f}{\partial x^2}\right)dt + \sigma \frac{\partial f}{\partial x}dW_t}%
    {Substituting the fundamental Itô multiplication rule $(dW_t)^2 = dt$ yields the deterministic drift correction term.}
  
  \vfill
  \BayesCaption{Blackboard Invariant: The non-zero quadratic variation of Brownian paths induces the Itô drift correction.}
\end{frame}

% =============================================================================
% 20. VISUAL ANALOGY: RECEPTIVE FIELDS
% =============================================================================
\begin{frame}{20 — Visual Analogy: Convolutional Receptive Fields}
  \begin{columns}[T,onlytextwidth]
    \begin{column}{.62\textwidth}
      \centering
      \begin{tikzpicture}[scale=.90,transform shape,
        label/.style={font=\scriptsize\bfseries\color{BayesSecondary},align=center}]
        % Image and explicitly highlighted local receptive field.
        \draw[BayesSecondary,thick] (0,0) grid[step=.48] (2.4,2.4);
        \fill[BayesTeal!20!white] (.96,.96) rectangle (1.92,1.92);
        \draw[BayesTeal,very thick] (.96,.96) rectangle (1.92,1.92);
        \node[label,below] at (1.2,-.12) {input image\\local $3\!\times\!3$ patch};
        \draw[bayes/flow] (2.65,1.44) -- (3.50,1.44);
        \node[bayes/accent,align=center,minimum width=1.20cm,minimum height=1.20cm,fill=BayesAmber] (kernel) at (4.25,1.44) {$3\!\times\!3$\\kernel};
        \draw[bayes/flow] (kernel.east) -- (5.35,1.44);
        \node[bayes/accent,align=center,minimum width=1.28cm,fill=BayesTeal] (feature) at (6.20,1.44) {feature\\activation};
        \draw[bayes/flow] (feature.south) -- ++(0,-.75);
        \node[label,below] at (6.20,.45) {deeper layers combine\\nearby activations};
      \end{tikzpicture}
    \end{column}
    \hfill
    \begin{column}{.34\textwidth}
      {\small\bfseries\color{BayesTeal}From local evidence to objects}\par\BayesGap
      {\small\textbf{Layer 1:} edges\par\BayesSmallGap
      \textbf{Layer 2:} textures\par\BayesSmallGap
      \textbf{Layer 3:} motifs\par\BayesSmallGap
      \textbf{Layer 4:} objects}\par\BayesGap
      \par\noindent\colorbox{BayesPale!30!white}{\parbox{.94\linewidth}{\centering\vspace{2.4mm}{\footnotesize\(\displaystyle y[i,j] = \sum_{u,v} x[i+u,j+v]w[u,v]\)}\vspace{2.4mm}}}
    \end{column}
  \end{columns}
  \vfill
  \BayesCaption{Visual explanation: a shared local kernel turns a receptive field into a feature activation; depth expands the effective field.}
\end{frame}

% =============================================================================
% 21. ENZYME KINETICS: REACTION COORDINATE
% =============================================================================
\begin{frame}{21 — Michaelis--Menten Enzyme Kinetics}
  \FigureDominant[.60]{%
    \begin{tikzpicture}
      \begin{axis}[
        bayes/science plot,
        xlabel={\small Substrate Concentration $[S]$ (mM)},
        ylabel={\small Reaction Rate $v$ ($\mu\text{M}/\text{s}$)},
        domain=0:10,
        samples=100,
        xmin=0, xmax=10, ymin=0, ymax=11.5,
        width=.98\linewidth, height=.54\textheight
      ]
        \addplot[very thick, color=BayesTeal] {10*x / (2 + x)};
        \addplot[thick, color=BayesSecondary, dashed] coordinates {(0,10) (10,10)};
        \node[above,font=\tiny\color{BayesSecondary}] at (axis cs:7,10) {Asymptote $V_{\max} = 10.0$};
        
        \addplot[thick, color=BayesCrimson, dotted] coordinates {(2,0) (2,5) (0,5)};
        \node[right,font=\tiny\color{BayesCrimson}] at (axis cs:2,5) {$(K_M, \frac{1}{2}V_{\max})$};
      \end{axis}
    \end{tikzpicture}%
  }{Hyperbolic Velocity Saturation}{%
    \begin{itemize}\setlength{\itemsep}{.25em}
      \item \textbf{First-Order Regime ($[S] \ll K_M$)}: Rate is linear in substrate $v \approx \frac{V_{\max}}{K_M}[S]$.
      \item \textbf{Zero-Order Regime ($[S] \gg K_M$)}: Active sites are fully saturated; rate reaches asymptote $V_{\max} = k_{\text{cat}}[E]_0$.
    \end{itemize}%
  }
  \vfill
  \BayesCaption{Scientific Reaction Coordinate: Direct visualization of asymptotic enzyme saturation.}
\end{frame}

% =============================================================================
% 22. QUANTUM STATE & BELL INEQUALITY PRECISION
% =============================================================================
\begin{frame}{22 — Quantum State \& Bell Non-Locality}
  \begin{columns}[c,onlytextwidth]
    \begin{column}{.48\textwidth}
      \centering
      \begin{quantikz}
        \lstick{$\ket{0}$} & \gate{H} & \ctrl{1} & \meter{} \\
        \lstick{$\ket{0}$} & \qw      & \targ{}  & \meter{}
      \end{quantikz}
      \vspace{.15cm}
      \[
        \ket{\Phi^+} = \frac{1}{\sqrt{2}}\left(\ket{00} + \ket{11}\right)
      \]
    \end{column}
    \hfill
    \begin{column}{.48\textwidth}
      {\small\bfseries\color{BayesTeal}Quantum Non-Local Correlations}\par\BayesGap
      {\small The entangled Bell state $\ket{\Phi^+}$ produces measurement correlations that violate the classical Bell--CHSH inequality.}\par\BayesGap
      \begin{SafeContentPanel}
        \centering
        $\displaystyle |S_{\mathrm{QM}}| = 2\sqrt{2} \approx 2.828 > 2$\par\BayesSmallGap
        {\footnotesize\color{BayesSecondary}(classical CHSH bound)}
      \end{SafeContentPanel}
      \BayesGap
      \BayesQuietNote{No local hidden-variable model reproduces these quantum predictions.}
    \end{column}
  \end{columns}
  \vfill
  \BayesCaption{Scientific Accuracy: Entanglement produces non-local correlations violating classical bounds without superluminal signalling.}
\end{frame}

% =============================================================================
% 23. FULLCANVAS SAFE-BOUNDING (Zero Edge Collision)
% =============================================================================
\begin{FullCanvas}
  \begin{tikzpicture}[remember picture,overlay]
    \fill[BayesCanvas] (current page.north west) rectangle (current page.south east);
    
    % Safe text boundary container (7.8mm margin)
    \node[anchor=north west,font=\sffamily\fontsize{16}{20}\selectfont\bfseries\color{BayesInk}]
      at ([xshift=0.85cm,yshift=-0.85cm]current page.north west) {23 — FullCanvas: Safe Edge-to-Edge Visualization};
    \node[anchor=north west,font=\small\color{BayesSecondary}]
      at ([xshift=0.85cm,yshift=-1.45cm]current page.north west) {Dedicated to panoramic multi-panel scientific graphics with zero header and footer chrome.};
    
    % Wave illustration
    \begin{scope}[shift={([xshift=2.2cm,yshift=-5.2cm]current page.north west)}]
      \draw[BayesTeal,very thick] (0,0) sin (2.4,1.4) cos (4.8,0) sin (7.2,-1.4) cos (9.6,0) sin (12.0,1.4);
      \draw[BayesBlue,very thick,dashed] (0,0) sin (2.4,0.7) cos (4.8,0) sin (7.2,-0.7) cos (9.6,0) sin (12.0,0.7);
      \draw[->,thick,color=BayesSlate] (-0.4,0) -- (12.4,0) node[right,font=\footnotesize] {$t$};
      \draw[->,thick,color=BayesSlate] (0,-1.8) -- (0,1.8) node[above,font=\footnotesize] {$\psi(t)$};
    \end{scope}
  \end{tikzpicture}
\end{FullCanvas}

% =============================================================================
% 24. DARKCANVAS: THEORETICAL PHASE TRANSITION
% =============================================================================
\begin{DarkCanvas}
  \begin{tikzpicture}[remember picture,overlay]
    \node[anchor=center,draw=BayesSecondary!55!BayesDark,line width=.45pt,rounded corners=1.2mm,
      fill=BayesDark,text width=.80\paperwidth,inner xsep=6mm,inner ysep=5mm,
      font=\BayesDisplayFont\fontsize{21}{26}\selectfont\bfseries\color{white},align=center]
      at (current page.center) {24 — Phase Transition:\\The Scaling Limit\\[4pt]
      {\normalsize\mdseries\color{BayesBorder}Asymptotic behavior when network width and depth approach infinity}};
  \end{tikzpicture}
\end{DarkCanvas}

% =============================================================================
% 25. HIGH-DENSITY KPI GRID
% =============================================================================
\begin{frame}{25 — Systems Performance Metrics}
  \StatGridFour{%
    \StatCard{3.12\,$\mu$s}{Kernel Latency}{CUDA XLA Fusion}{Zero host-device dispatch overhead}%
  }{%
    \StatCard{94.2\%}{GPU Utilization}{Tensor Core Saturation}{Full matrix-multiply occupancy}%
  }{%
    \StatCard{1.8\,$\times$}{KV Cache Saving}{PagedAttention vLLM}{Dynamic non-contiguous memory}%
  }{%
    \StatCard{0.001}{Error Tolerance}{FP16 Mixed Precision}{No loss of downstream accuracy}%
  }
  \vfill
  \TakeawayBanner{System throughput scales linearly with batch dimension when memory bandwidth is fully saturated.}
\end{frame}

% =============================================================================
% 26. MULTI-COLUMN TAXONOMY
% =============================================================================
\begin{frame}{26 — Probabilistic Model Taxonomy}
  \begin{columns}[T,onlytextwidth]
    \begin{column}{.31\textwidth}
      {\small\bfseries\color{BayesTeal}Analytically Tractable}\par\BayesSmallGap
      {\footnotesize
        \begin{itemize}\setlength{\itemsep}{.2em}
          \item Conjugate Exponential Families
          \item Kalman Filter (linear-Gaussian)
          \item GP Regression (Gaussian likelihood)
        \end{itemize}
      }
      \BayesSmallGap
      {\scriptsize\color{BayesMuted}Exact only under the stated model assumptions.}
    \end{column}
    \hfill{\color{BayesBorder}\vrule width .4pt}\hfill
    \begin{column}{.31\textwidth}
      {\small\bfseries\color{BayesBlue}Asymptotically Exact MCMC}\par\BayesSmallGap
      {\footnotesize
        \begin{itemize}\setlength{\itemsep}{.2em}
          \item Metropolis-Hastings
          \item Hamiltonian Monte Carlo (HMC)
          \item No-U-Turn Sampler (NUTS)
        \end{itemize}
      }
      \BayesSmallGap
      {\scriptsize\color{BayesMuted}Practicality depends on geometry, mixing, and available compute.}
    \end{column}
    \hfill{\color{BayesBorder}\vrule width .4pt}\hfill
    \begin{column}{.31\textwidth}
      {\small\bfseries\color{BayesAmber}Deterministic Approximation}\par\BayesSmallGap
      {\footnotesize
        \begin{itemize}\setlength{\itemsep}{.2em}
          \item Variational Inference (VI)
          \item Expectation Propagation
          \item Laplace Approximation
        \end{itemize}
      }
      \BayesSmallGap
      {\scriptsize\color{BayesMuted}Can exploit minibatches and structure; quality needs diagnosis.}
    \end{column}
  \end{columns}
  \vfill
  \BayesCaption{Taxonomy: Structured classification across precision, computational scaling, and asymptotic guarantees.}
\end{frame}

% =============================================================================
% 27. SIDE-BY-SIDE DUAL IMPLEMENTATION COMPARISON
% =============================================================================
\begin{frame}[fragile]{27 — Vectorized vs Iterative Execution}
  \begin{columns}[T,onlytextwidth]
    \begin{column}{.48\textwidth}
      {\small\bfseries\color{BayesCrimson}Anti-Pattern: Python Iterative Loop}\par\vspace{.04cm}
\begin{CodeBlock}{Python}
# Slow: Interpreted Python loop
res = []
for i in range(len(x)):
    res.append(x[i] * w[i])
\end{CodeBlock}
      \vspace{.04cm}
      {\scriptsize\color{BayesSecondary}\textbf{Latency:} 42.1\,ms (GIL overhead)}
    \end{column}
    \hfill{\color{BayesBorder}\vrule width .4pt}\hfill
    \begin{column}{.48\textwidth}
      {\small\bfseries\color{BayesTeal}Optimal: SIMD Vectorized Tensor}\par\vspace{.04cm}
\begin{CodeBlock}{Python}
# Fast: Fused vector operation
res = np.dot(x, w)
\end{CodeBlock}
      \vspace{.04cm}
      {\scriptsize\color{BayesSecondary}\textbf{Latency:} 0.08\,ms (520$\times$ speedup via AVX-512)}
    \end{column}
  \end{columns}
  \vfill
  \TakeawayBanner{Vectorized SIMD memory layouts eliminate interpreter dispatch overhead.}
\end{frame}

% =============================================================================
% 28. PRODUCTION CASE STUDY & ROOT CAUSE ANALYSIS
% =============================================================================
\begin{frame}{28 — Production Case Study: Deadlock Under Contention}
  \begin{columns}[T,onlytextwidth]
    \begin{column}{.48\textwidth}
      {\small\bfseries\color{BayesCrimson}Failure Mode: Distributed Deadlock}\par\vspace{.06cm}
      {\small Under peak burst load, asynchronous workers acquired tensor mutex locks in non-deterministic order: Worker 1 locked $T_A$ and requested $T_B$, while Worker 2 locked $T_B$ and requested $T_A$, triggering cyclic lock freeze.}\par\vspace{.06cm}
      {\scriptsize\color{BayesSecondary}Resulted in 100\% worker stall and training pipeline freeze.}
    \end{column}
    \hfill{\color{BayesBorder}\vrule width .4pt}\hfill
    \begin{column}{.48\textwidth}
      {\small\bfseries\color{BayesTeal}Production Fix: Canonical Lock Ordering}\par\vspace{.06cm}
      {\small Enforced global strict lexicographical lock ordering on all tensor keys $\text{ID}(T_i) < \text{ID}(T_j)$ prior to acquisition.}\par\vspace{.06cm}
      {\scriptsize\color{BayesSecondary}Mathematically eliminates cyclic wait conditions (Dijkstra's resource hierarchy theorem).}
    \end{column}
  \end{columns}
  \vspace{.20cm}
  \DataProvenance[INCIDENT POST-MORTEM]{Cluster Telemetry Severity-1 Resolution}
\end{frame}

% =============================================================================
% 29. PROGRESSIVE REVEAL / COGNITIVE STATE SEQUENCE
% =============================================================================
\begin{frame}{29 — Progressive Cognitive Sequence}
  \begin{center}
    {\small\bfseries\color{BayesTeal}The 3-Step Cognitive Revelation of Variational Inference}\par\vspace{.25cm}
  \end{center}
  \begin{columns}[T,onlytextwidth]
    \begin{column}{.31\textwidth}
      {\small\bfseries Step 1: Intractable Target}\par\vspace{.06cm}
      \[
        p(\theta \mid \mathcal{D}) = \frac{p(\mathcal{D}, \theta)}{p(\mathcal{D})}
      \]
      {\scriptsize The true posterior evidence integral $p(\mathcal{D})$ cannot be evaluated in high dimensions.}
    \end{column}
    \hfill
    \begin{column}{.31\textwidth}
      {\small\bfseries Step 2: Variational Family}\par\vspace{.06cm}
      \[
        q_\phi(\theta) \approx p(\theta \mid \mathcal{D})
      \]
      {\scriptsize Posit a tractable parametric density family $q_\phi$ (e.g. Gaussian mean-field).}
    \end{column}
    \hfill
    \begin{column}{.31\textwidth}
      {\small\bfseries Step 3: ELBO Optimization}\par\vspace{.06cm}
      \[
        \arg\max_\phi \mathbb{E}_{q_\phi}[\ln p(\mathcal{D},\theta)] - \mathrm{KL}
      \]
      {\scriptsize Inference converts into standard numerical optimization via gradient ascent.}
    \end{column}
  \end{columns}
  \vfill
  \TakeawayBanner{Inference becomes optimization: trade asymptotic exactness for deterministic scalability.}
\end{frame}

% =============================================================================
% 30. SCHOLARLY REFERENCES & CURRICULUM SYNTHESIS
% =============================================================================
\begin{References}
  \item \textbf{Gelman, A., et al.} (2013). \emph{Bayesian Data Analysis} (3rd ed.). Chapman and Hall/CRC.
  \item \textbf{Jaynes, E. T.} (2003). \emph{Probability Theory: The Logic of Science}. Cambridge University Press.
  \item \textbf{Hoffman, M. D., \& Gelman, A.} (2014). The No-U-Turn Sampler. \emph{J. Mach. Learn. Res.}, 15(1), 1593--1623.
  \item \textbf{Vaswani, A., et al.} (2017). Attention Is All You Need. \emph{NeurIPS}, 30, 5998--6008.
\end{References}

\end{document}
